% BHU PYQ -- Manually transcribed & error-corrected
\documentclass[11pt,a4paper]{article}

\usepackage[a4paper,top=2.5cm,bottom=2.5cm,left=2.8cm,right=2.8cm,headheight=14pt]{geometry}
\usepackage{amsmath,amssymb}
\usepackage[T1]{fontenc}
\usepackage[utf8]{inputenc}
\usepackage{lmodern}
\usepackage{microtype}
\usepackage{fancyhdr}
\usepackage{enumitem}
\usepackage{hyperref}
\usepackage{lastpage}
\usepackage{parskip}

\hypersetup{colorlinks=true,urlcolor=black,linkcolor=black,
  pdftitle={BPE-602: Elementary Nanoscience and Exotic Materials -- BHU PYQ},pdfauthor={Banaras Hindu University}}

\pagestyle{fancy}\fancyhf{}
\renewcommand{\headrulewidth}{0.4pt}\renewcommand{\footrulewidth}{0.4pt}
\fancyhead[L]{\small   \textit{Banaras Hindu University}}
\fancyhead[R]{\small   \textit{BPE-602: Elementary Nanoscience and Exotic Materials}}
\fancyfoot[C]{\small Page  \thepage\ of \pageref{LastPage}}


\newlist{parts}{enumerate}{2}
\setlist[parts,1]{label=(\alph*),leftmargin=*,itemsep=5pt,topsep=3pt}
\setlist[parts,2]{label=(\roman*),leftmargin=*,itemsep=3pt,topsep=2pt}

\newcommand{\pts}[1]{\hfill   \textit{[#1]}}


\begin{document}


\begin{center}
  {\large\bfseries BANARAS HINDU UNIVERSITY}\\[2pt]
  {
\normalsize B.Sc. (Hons.) Semester VI Examination, 2022-23}\\[6pt]
  \rule{0.8\linewidth}{0.5pt}\\[6pt]
  {\large\bfseries Paper No. BPE-602: Elementary Nanoscience and Exotic Materials}\\[4pt]
  {
\normalsize Time: 3 Hours\hfill Maximum Marks: 70}
\end{center}
\rule{\linewidth}{0.4pt}
\smallskip



\noindent   \textit{Instructions: Answer all questions. Symbols have their usual meanings.}
\bigskip




\noindent   \textbf{Question 1.} Answer in brief any   \textbf{four} of the following: \pts{4 $ \times$ 3.5 = 14}

\begin{parts}
  \item Explain how the band gap of a material changes with a change in its dimensions.
  \item How is the crystallite size determined using X-ray diffraction (XRD) patterns?
  \item Explain the Fibonacci sequence.
  \item What is meant by surface plasmon resonance? How is it related to nanomaterials?
  \item What are Type-I and Type-II superconductors?
\end{parts}

\medskip\rule{\linewidth}{0.2pt}




\noindent   \textbf{Question 2.}

\begin{parts}
  \item What are the properties of high-$T_c$ superconductors? Explain how they are different from conventional superconductors. \pts{8}
  \item What is the Meissner effect? A superconducting lead has a critical temperature of $7.26\, \text{K}$ at zero magnetic field and a critical field of $8  \times 10^5\, \text{A/m}$ at $0\, \text{K}$. Find the critical field at $5\, \text{K}$. \pts{6}
\end{parts}

\medskip\rule{\linewidth}{0.2pt}




\noindent   \textbf{Question 3.}

\begin{parts}
  \item What are colossal magnetoresistance (CMR) materials? Discuss the double exchange mechanism of colossal magnetoresistance. \pts{8}
  \item How is CMR different from ordinary resistance and Giant Magnetoresistance (GMR)? \pts{6}
\end{parts}
\hfill   \textbf{OR}

\begin{parts}
  \item What are carbon nanotubes (CNTs)? Describe the armchair, zigzag, and chiral types of CNTs using suitable diagrams. \pts{8}
  \item Write in brief the working of a transmission electron microscope (TEM). \pts{6}
\end{parts}

\medskip\rule{\linewidth}{0.2pt}




\noindent   \textbf{Question 4.}

\begin{parts}
  \item Explain the particle size effect on the optical spectra of nanomaterials. Explain how the optical spectra can be used to determine the size of nanoparticles. \pts{9}
  \item Explain how the band gap of nanoparticles is evaluated using their UV-Vis absorption spectra. \pts{5}
\end{parts}

\medskip\rule{\linewidth}{0.2pt}




\noindent   \textbf{Question 5.}

\begin{parts}
  \item What are quasicrystals? Give one example. Provide the construction of a 1D quasicrystal by projecting a 2D periodic lattice onto a line with an irrational slope. \pts{8}
  \item What is a Quantum Dot (QD)? If the size of a QD is halved, what will happen to the frequency of the fluorescent light emitted by it? \pts{6}
\end{parts}
\hfill   \textbf{OR}

\begin{parts}
  \item Describe the differences between the physical and chemical synthesis routes of nanomaterials. Explain the working principle and limitations of high-energy ball milling. \pts{9}
  \item Describe the bio-route of synthesis of nanoparticles. \pts{5}
\end{parts}

\vfill
\rule{\linewidth}{0.4pt}

\begin{center}\small
  \href{https://bhu-exam.vercel.app/}{Click here to visit our website}\\[4pt]
  \href{https://me-aryan.vercel.app/}{Click here to visit our contributor}\\[4pt]
  \href{https://quashan-me.vercel.app/}{Click here to visit our contributor}
\end{center}

\end{document}
